%------------------------------------------------------------------------------%
\begin{question}
  Determine a posição relativa entre a reta e o plano \\
  e o ponto de interseção, se ele existir
  \[
    r\colon
    \begin{cases}
      x = 1 + t \\
      y = 2 - t \\
      z = 3
    \end{cases}
  \]
  \[
    \gamma\colon x + y - 3 = 0
  \]

\end{question}

%------------------------------------------------------------------------------%
\begin{resolution}{Posição Relativa}
  \onslide<+->{O vetor diretor da reta é \qquad
  \(
    \vec{v} = (1, -1, 0)
  \)}

  \onslide<+->{O vetor normal do plano é \quad
  \(
    \vec{n} = (1, 1, 0)
  \)}
  \begin{align*}
    \onslide<+->{\vec{v} \cdot \vec{n}}
    \onslide<+->{& = 1 \times 1 + (- 1)1 + 0 \times 0} \\
    \onslide<+->{& = 1 - 1 + 0} \\
    \onslide<+->{& = 0}
  \end{align*}
  \onslide<+->{a reta é paralela ao plano ou está contida nele}
\end{resolution}

%------------------------------------------------------------------------------%
\begin{resolution}{Ponto de Interseção}
  Escolhemos o ponto da reta
  \[
    A = (1, 2, 3)
  \]
  \pause
  e substituimos na equação do plano
  \pause
  \[
    x + y - 3
    \pause
    = 1 + 2 - 3
    \pause
    = 0
  \]
  \pause
  Como o ponto da reta satisfaz a equação do plano
  \[
    \pause
    A\in\gamma
  \]
\end{resolution}

%------------------------------------------------------------------------------%
\begin{resolution}{Ponto de Interseção}
  Como a reta possui um ponto no plano e é paralela ao plano

  \pause
  todos os seus pontos pertencem ao plano
  \pause
  \[
    r\subset\gamma
  \]
  A reta está \emph{contida no plano}
\end{resolution}

%------------------------------------------------------------------------------%
\endinput
